\documentclass[12pt,a4paper]{article}
\usepackage[utf8]{inputenc}
\usepackage[T1]{fontenc}
\usepackage[polish]{babel}
\usepackage{geometry}
\usepackage{hyperref}
\usepackage{listings}
\usepackage{xcolor}
\usepackage{booktabs}
\usepackage{enumitem}
\usepackage{parskip}

\geometry{margin=2.5cm}

\hypersetup{
    colorlinks=true,
    linkcolor=blue,
    urlcolor=blue
}

\lstset{
    basicstyle=\ttfamily\small,
    breaklines=true,
    frame=single,
    backgroundcolor=\color{gray!10},
    keywordstyle=\color{blue},
    commentstyle=\color{green!60!black},
    stringstyle=\color{red!70!black}
}

\title{\textbf{Sztuczna Inteligencja i Inżynieria Wiedzy}\\
\large Zadanie 1 -- Raport Implementacyjny}
\author{}
\date{2026}

\begin{document}

\maketitle
\tableofcontents
\newpage

% ============================================================
\section{Struktura projektu}
% ============================================================

Projekt realizuje wyszukiwanie połączeń kolejowych na danych GTFS Kolei Dolnośląskich.
Kod źródłowy zorganizowany jest w pakiecie \texttt{src/} z podziałem na moduły:

\begin{lstlisting}
src/
  config.py              # stale konfiguracyjne
  algorithms/
    astar.py             # algorytm A*
    dijkstra.py          # algorytm Dijkstry
    tabu_search.py       # Tabu Search (warianty 2a-2d)
    common.py            # wspolne funkcje (relaksacja krawedzi, rekonstrukcja sciezki)
  gtfs/
    loader.py            # wczytywanie plikow GTFS
    graph.py             # budowanie grafu polaczen
    calendar.py          # filtrowanie kursow wg daty
  cli/
    pathfinder_cli.py    # interfejs CLI -- Zadanie 1
    tsp_cli.py           # interfejs CLI -- Zadanie 2
  utils/
    time_utils.py        # konwersje czasu
    geo.py               # odleglosc haversine
server.py                # serwer HTTP z API REST
gtfs_visualizer.html     # interfejs webowy
\end{lstlisting}

% ============================================================
\section{Zadanie 1 -- Wyszukiwanie najkrótszej ścieżki}
% ============================================================

\subsection{Algorytm Dijkstry (1a)}

\textbf{Plik:} \texttt{src/algorithms/dijkstra.py}\\
\textbf{Funkcja:} \texttt{dijkstra(graph, starts, ends, criterion)}

Implementacja klasycznego algorytmu Dijkstry z kolejką priorytetową (\texttt{heapq}).
Obsługuje oba kryteria optymalizacji: czas przejazdu (\texttt{t}) oraz liczbę przesiadek
(\texttt{p}).

Uruchamiany z poziomu CLI flagą \texttt{--dijkstra}:
\begin{lstlisting}[language=bash]
python -m src.cli.pathfinder_cli "Wroclaw Glowny" "Jelenia Gora" t \
    "2026-03-09 08:00" --dijkstra
\end{lstlisting}

\subsection{Algorytm A* -- kryterium czasu (1b)}

\textbf{Plik:} \texttt{src/algorithms/astar.py}\\
\textbf{Funkcja główna:} \texttt{astar(graph, starts, ends, criterion, ...)}\\
\textbf{Heurystyka:} \texttt{heuristic\_time(node, end\_stops, gtfs\_data)}

Heurystyka czasowa szacuje pozostały koszt jako odległość geograficzną (wzór haversine,
funkcja \texttt{distance\_km} w \texttt{src/utils/geo.py}) podzieloną przez maksymalną
prędkość pociągu (\texttt{MAX\_SPEED\_KMH} w \texttt{src/config.py}).
Gwarantuje to dopuszczalność heurystyki -- nigdy nie przeszacowuje rzeczywistego czasu.

\subsection{Algorytm A* -- kryterium przesiadek (1c)}

\textbf{Plik:} \texttt{src/algorithms/astar.py}\\
\textbf{Heurystyka:} \texttt{heuristic\_transfers(node, end\_stops, direct\_set)}\\
\textbf{Funkcja pomocnicza:} \texttt{build\_direct\_stop\_set(graph, end\_stops)}

Heurystyka przesiadkowa zwraca 0, jeśli istnieje bezpośrednie połączenie z bieżącego
węzła do stacji docelowej (precomputed \texttt{direct\_set}), a w przeciwnym razie
zwraca 1. Jest to heurystyka dopuszczalna -- nigdy nie zawyża liczby przesiadek.

\subsection{Modyfikacje przyspieszające obliczenia (1d)}

\textbf{Plik:} \texttt{src/algorithms/astar.py}

Zaimplementowane usprawnienia:
\begin{itemize}
    \item \textbf{Wczesne zatrzymanie} (\textit{early exit}): algorytm kończy działanie
          natychmiast po wyciągnięciu węzła docelowego z kolejki, bez eksploracji
          pozostałych gałęzi.
    \item \textbf{Grupowanie przystanków stacji} (\texttt{station\_group} w
          \texttt{src/gtfs/graph.py}): zamiast jednego węzła startowego, do kolejki
          dodawane są wszystkie perony danej stacji, co eliminuje zbędne krawędzie
          transferowe na początku trasy.
    \item \textbf{Pominięcie zdominowanych stanów}: jeśli węzeł zostaje wyciągnięty
          z kolejki z kosztem wyższym niż dotychczas znane optimum dla danego przystanku,
          jest pomijany (\texttt{if cost > dist[node]: continue}).
    \item \textbf{Precomputed \texttt{direct\_set}} dla heurystyki przesiadkowej:
          jednorazowo obliczany w \texttt{build\_direct\_stop\_set()} i przekazywany
          jako parametr, zamiast obliczać przy każdym wywołaniu heurystyki.
\end{itemize}

\subsection{Interfejs CLI -- Zadanie 1}

\textbf{Plik:} \texttt{src/cli/pathfinder\_cli.py}\\
\textbf{Funkcja:} \texttt{main()} / \texttt{\_run()}\\
\textbf{Formatowanie wyników:} \texttt{print\_results()}

\begin{lstlisting}[language=bash]
# Przyklad uzycia (A*, kryterium czasu):
python -m src.cli.pathfinder_cli "Wroclaw Glowny" "Legnica" t "2026-03-10 07:30"

# Dijkstra, kryterium przesiadek:
python -m src.cli.pathfinder_cli "Swidnica Miasto" "Walbrzych Glowny" p \
    "2026-03-10 09:00" --dijkstra
\end{lstlisting}

Na standardowe wyjście wypisywane są kolejne odcinki trasy (przystanek początkowy,
końcowy, linia, czas odjazdu/przyjazdu). Na standardowe wyjście błędów -- wartość
zminimalizowanego kryterium i czas obliczeń.

% ============================================================
\section{Zadanie 2 -- Problem komiwojażera (Tabu Search)}
% ============================================================

\subsection{Interfejs CLI -- Zadanie 2}

\textbf{Plik:} \texttt{src/cli/tsp\_cli.py}\\
\textbf{Funkcja:} \texttt{main()}

\begin{lstlisting}[language=bash]
# Przyklad uzycia (Tabu Search z aspiracja i probkowaniem):
python -m src.cli.tsp_cli "Wroclaw Glowny" "Legnica;Jelenia Gora;Swidnica" \
    t "2026-03-10 08:00" --2c --2d
\end{lstlisting}

Flagi wyboru wariantu Tabu Search: \texttt{--2a}, \texttt{--2b} (domyślnie),
\texttt{--2c} (aspiracja), \texttt{--2d} (próbkowanie sąsiedztwa).

\subsubsection*{Macierz kosztów}

\textbf{Funkcja:} \texttt{compute\_cost\_matrix()} -- oblicza macierz kosztów
przejazdu między wszystkimi parami przystanków za pomocą A* (optymistyczna wersja
bez uwzględniania czasu przyjazdu).\\
\textbf{Funkcja:} \texttt{compute\_cost\_matrix\_chained()} -- iteracyjne przeliczanie
kosztów z rzeczywistymi czasami przyjazdów (udoskonalenie wyniku Tabu Search).

\subsection{Tabu Search bez ograniczenia T (2a)}

\textbf{Plik:} \texttt{src/algorithms/tabu\_search.py}\\
\textbf{Funkcja:} \texttt{tabu\_search\_unlimited(cost\_matrix, n, criterion)}

Podstawowy algorytm Tabu Search zgodny z algorytmem Knoxa. Lista tabu \texttt{T}
rośnie nieograniczenie przez cały czas działania. Ruchem jest zamiana dwóch
przystanków w trasie (\textit{2-opt reversal}).

\subsection{Tabu Search ze zmiennym rozmiarem T (2b)}

\textbf{Plik:} \texttt{src/algorithms/tabu\_search.py}\\
\textbf{Funkcja:} \texttt{tabu\_search\_variable\_tenure(cost\_matrix, n, criterion)}

Rozmiar kolejki FIFO tablicy tabu wynosi $t = \lceil |L| / 2 \rceil$, gdzie $|L|$
to liczba przystanków pośrednich. Dla małych tras lista jest krótka (mniejsze
zużycie pamięci), dla długich -- dłuższa (lepsza ochrona przed cyklami).

\subsection{Kryterium aspiracji (2c)}

\textbf{Plik:} \texttt{src/algorithms/tabu\_search.py}\\
\textbf{Funkcja:} \texttt{tabu\_search\_core(..., use\_aspiration=True)}

Aktywowane flagą \texttt{use\_aspiration=True} w \texttt{tabu\_search\_core()}.
Kryterium aspiracji: jeśli ruch jest na liście tabu, ale prowadzi do rozwiązania
\emph{lepszego} niż dotychczasowe najlepsze \texttt{best\_cost}, jest dopuszczany.

\subsection{Próbkowanie sąsiedztwa (2d)}

\textbf{Plik:} \texttt{src/algorithms/tabu\_search.py}\\
\textbf{Funkcja:} \texttt{tabu\_search\_core(..., use\_sampling=True)}

Aktywowane flagą \texttt{use\_sampling=True}. Zamiast przeszukiwać wszystkie
$O(n^2)$ sąsiadów 2-opt, losowo próbkowana jest podzbiorowa liczba kandydatów
(\texttt{TABU\_SAMPLE\_SIZE} z \texttt{src/config.py}). Skraca czas iteracji
kosztem ewentualnej gorszej jakości rozwiązania.

\subsection{Inicjalizacja trasy i funkcja kosztu}

\textbf{Plik:} \texttt{src/algorithms/tabu\_search.py}\\
\textbf{Funkcje:}
\begin{itemize}
    \item \texttt{nearest\_neighbour\_init()} -- zachłanna inicjalizacja trasą
          metodą najbliższego sąsiada,
    \item \texttt{random\_init()} -- losowa permutacja przystanków,
    \item \texttt{tour\_cost\_matrix()} -- oblicza całkowity koszt trasy na
          podstawie macierzy kosztów.
\end{itemize}

\subsection{Ewaluacja trasy}

\textbf{Plik:} \texttt{src/cli/tsp\_cli.py}\\
\textbf{Funkcja:} \texttt{evaluate\_tour()}

Po znalezieniu optymalnej kolejności przystanków, \texttt{evaluate\_tour()} symuluje
rzeczywistą podróż z uwzględnieniem rzeczywistych czasów odjazdu/przyjazdu, kursów
oraz przesiadek. Obsługuje też bezpośrednie przejazdy (\texttt{check\_direct\_connections},
\texttt{find\_continuation}).

% ============================================================
\section{Wczytywanie danych GTFS i budowanie grafu}
% ============================================================

\textbf{Plik:} \texttt{src/gtfs/loader.py}\\
\textbf{Funkcja:} \texttt{load\_gtfs(path)} -- wczytuje wszystkie pliki GTFS do
słowników Pythona.

\textbf{Plik:} \texttt{src/gtfs/graph.py}\\
\textbf{Funkcja:} \texttt{build\_graph(gtfs\_data, start\_date, num\_days, ...)} --
buduje skierowany graf zależny od czasu:
\begin{itemize}
    \item krawędzie kursów ze \texttt{stop\_times.txt} (dla aktywnych serwisów),
    \item krawędzie transferowe między przystankami tej samej stacji
          (\texttt{\_add\_transfer\_edges()}).
\end{itemize}

\textbf{Plik:} \texttt{src/gtfs/calendar.py}\\
\textbf{Funkcja:} \texttt{is\_service\_active(service\_id, date, ...)} -- sprawdza,
czy dany kurs jest aktywny w podanej dacie, uwzględniając wyjątki z
\texttt{calendar\_dates.txt}.

Wyszukiwanie przystanków po nazwie: \texttt{find\_stop()} w \texttt{src/gtfs/graph.py}
-- dopasowanie dokładne, z sufiksem,,główny'', prefiksowe, a następnie podciąg nazwy.

% ============================================================
\section{Biblioteki zewnętrzne}
% ============================================================

\begin{itemize}
    \item \textbf{Python 3.x} -- język implementacji.
    \item \textbf{heapq} (biblioteka standardowa) -- kolejka priorytetowa w
          algorytmach Dijkstry i A*.
    \item \textbf{csv} (biblioteka standardowa) -- wczytywanie plików GTFS.
    \item \textbf{datetime} (biblioteka standardowa) -- obsługa dat i kalendarza kursów.
    \item \textbf{math} (biblioteka standardowa) -- obliczenia trygonometryczne
          (haversine).
    \item \textbf{random} (biblioteka standardowa) -- losowe próbkowanie sąsiedztwa
          w Tabu Search.
    \item \textbf{http.server} (biblioteka standardowa) -- serwer HTTP
          (\texttt{server.py}).
    \item \textbf{itertools} (biblioteka standardowa) -- generowanie permutacji
          (brute force dla $n \le 4$ w \texttt{tsp\_cli.py}).
\end{itemize}

Projekt nie korzysta z zewnętrznych bibliotek spoza biblioteki standardowej Pythona.

% ============================================================
\section{Napotkane problemy}
% ============================================================

\begin{itemize}
    \item \textbf{Czasy przekraczające 24h w GTFS} -- plik \texttt{stop\_times.txt}
          zawiera czasy w formacie \texttt{HH:MM:SS}, gdzie \texttt{HH} może
          przekraczać 23 (np. \texttt{25:10:00} oznacza 1:10 następnego dnia).
          Rozwiązano to przez konwersję wszystkich czasów do sekund od północy
          (\texttt{time\_to\_sec()} w \texttt{src/utils/time\_utils.py}) i obsługę
          wartości $> 86400$.

    \item \textbf{Hierarchia stacji w GTFS} -- przystanki mogą być peronami
          (podrzędnymi) lub stacjami (nadrzędnymi). Krawędzie transferowe muszą
          łączyć \emph{wszystkie} perony tej samej stacji, nie tylko pary o
          identycznym \texttt{stop\_id}. Obsłużono to w \texttt{\_add\_transfer\_edges()}.

    \item \textbf{Dopuszczalność heurystyki A*} -- heurystyka musi być dopuszczalna
          (nie przeszacowywać). Dla kryterium czasu użyto odległości geograficznej
          podzielonej przez \texttt{MAX\_SPEED\_KMH}. Dla kryterium przesiadek wartość
          0/1 jest zawsze dopuszczalna.

    \item \textbf{Macierz kosztów a rzeczywiste czasy} -- optymistyczna macierz
          kosztów (A* bez uwzględniania łańcucha czasu) może być nieadekwatna po
          ustaleniu kolejności przystanków. Rozwiązano to przez iteracyjne przeliczanie
          macierzy z rzeczywistymi czasami przyjazdów (\texttt{compute\_cost\_matrix\_chained()}).

    \item \textbf{Identyfikacja przesiadek} -- zmiana \texttt{trip\_id} nie zawsze
          oznacza przesiadkę (ten sam pojazd może kontynuować kurs pod innym
          identyfikatorem \texttt{block\_id}). Obsłużono to w \texttt{find\_continuation()}
          i \texttt{check\_direct\_connections()}.
\end{itemize}

\end{document}
